\section{Limitations and Conclusion}

The evaluation uses synthetic tools and attack scenarios, a limited set of malicious scenario classes, receipt-based correlation in the evaluated path, one reference gateway, local harness performance measurements, and one frozen model/runtime selection. Strict action-trace replay approximates adaptive behavior after denial; it does not model an agent that changes its strategy after each blocked call. TCOP is not evaluated on a production incident, Internet-scale federation, or real-world attack frequency.

TCOP is useful only when a participating observer detects a useful fact in time and the receiver has an enforceable local boundary. It cannot create detections, guarantee observation truth, prevent an undetected origin compromise, or reverse completed harm. The evaluated central comparison is bounded and does not support universal claims about centralized defense.

Within those limits, the evidence supports a systems-security conclusion: signed bounded context can reduce downstream blast radius when it becomes actionable inside a containment window, while preserving receiver-local decision authority. The accompanying artifact makes the frozen deterministic and replay results reproducible without live-model credentials and retains the observed availability cost as part of the design trade-off.

\paragraph{Future work: an independent sandbox runtime as the receiver boundary.}
The physical reference path in Section~\ref{sec:agent-gateway} enforces Domain-B decisions through a reference gateway that we built, and the two \texttt{tcopd} endpoints ran as separate processes rather than on separately administered hosts. The pre-registered cross-host study (Appendix~\ref{app:amendments}) already requires two distinct non-local hosts and remains sealed as a blocked artifact. We intend to execute it with the receiver boundary supplied by an independently maintained agent sandbox runtime rather than our reference gateway. Out-of-process runtimes such as NVIDIA OpenShell, as packaged in NemoClaw, expose a declarative policy over filesystem, process, network egress, and inference routing, with deny-by-default egress and an operator approval path for policy change. Such a runtime is a receiver-local enforcement point in the sense this paper requires: it is outside the agent process, it does not depend on the agent's cooperation, and it already distinguishes allowed destinations from unknown ones. TCOP's role in that configuration is narrow and testable. A validated TCX may only tighten the runtime policy (remove an egress destination, revoke a credential scope, freeze or snapshot the sandbox); it may never widen it, and it never issues the enforcement call itself. The planned arm places two live agents on two hosts in declared adversarial roles, with Domain A observing and signing, Domain B deciding through the runtime policy, and the same five fixed adaptive branches of Amendment~002 driving the attacker. The July 2026 Hugging Face intrusion~\cite{hf-timeline-2026,metr-hf-2026} is the motivating case: every cross-domain step in that incident was an allowed network call from the agent's own side, which a per-action approval gate cannot distinguish from benign use and a receiver-side egress policy can, provided the observing domain's evidence reaches it in time. The claim boundary is unchanged: where receiver-bound evidence would have had a role, not that it would have stopped the incident.

\begin{table*}[t]\centering\small
\caption{What the evaluation supports and does not support.}
\begin{tabular}{p{0.3\textwidth}p{0.6\textwidth}}\toprule
Evidence & Supported scope \\ \midrule
Deterministic matched pairs & Causal comparison under matched frozen local policy and synthetic scenarios \\
Executable replay & Exact captured action-sequence comparison without model credentials \\
Live trace generation & A bounded credentialed cohort generated action attempts \\
Single-environment reference-path deployment & A signed tcopd-A-to-tcopd-B path and receiver-local gateway decisions in one reference environment \\
Validation-value and assurance cohort & Semantic binding, hostile-peer handling, receipt correlation, and receiver-validator behavior in a credential-free synthetic cohort \\
Post-denial substitution replay & Policy-conditioned post-denial branches in a bounded reference runtime, not deployment prevalence \\
Campaign Correlation frontier & Deterministic receiver-local Campaign Correlation selectivity, not a general escalation policy or deployment-effectiveness estimate \\
Full external-warning ledger & All 168 sealed external source cases: 25 accepted exact bindings, 47 attack-bearing no-warning cases, and 96 benign no-warning cases; authorization coverage, not detector-quality evidence \\
Not established & Production deployment, Internet-scale federation, key-lifecycle operation, uninstrumented incidents, or universal superiority to central defense \\ \bottomrule
\end{tabular}\label{tab:limitations}\end{table*}

